\label{sec:solution}

From a user's perspective, the whole pipeline is launched with a single orchestration script that opens a multi-window terminal session and starts, in order: the middleware router required for reliable playback of the large camera recordings, the camera and navigation bag playback (kept time-synchronized), the pose-broadcasting node, the static sensor-frame transforms, a 3D visualization of the sensor rig, the visualization tool itself, the rectification node, the chosen disparity backend, the point-cloud reconstruction node, the altimeter-publishing node, and the two wave-parameter estimators. A single configuration switch selects which of the six disparity backends to run, making the quality/speed/hardware trade-off discussed in Section~\ref{sec:architecture} a one-line choice rather than a code change.

Once running, the visualization tool shows, in a common georeferenced frame: the live colourized 3D point cloud of the sea surface, the platform's estimated pose and travelled path, the raw lidar point cloud, the rig's sensor frames, and the live disparity map. The two wave-estimation nodes publish their sea-state parameters (height, period, and diagnostic quality figures) as ROS~2 topics, which can be inspected live or logged for offline analysis; this is also how the cross-validation and comparison numbers in Section~\ref{sec:validation} were obtained.

\begin{figure}[h]
    \centering
    \includegraphics[width=\textwidth]{full_rviz}
    \caption{RViz visualization of the running pipeline. The 3D viewport (right) shows the georeferenced sea-surface point cloud and the sensor rig wireframe, with the platform's travelled path (green) overlaid in the \texttt{map} (ENU) frame. The left panel shows live image feeds: the rainbow-coded disparity map, the rectified left camera image, the thermal camera, and the horizon-detection output with the sky-exclusion line marked in red.}
    \label{fig:rviz-screenshot}
\end{figure}

The sensor frame layout used by the visualization is shown in Figure~\ref{fig:tf-frame-tree}.
